\chapter{[RWT] Rettungswesen und Einsatztaktik}

\section{Kommunikationswege}

\subsection{Funk}

\section{Einsatzmittel}

Im deutschsprachigen Ram hat sich die Einteilung in die
Typen 

\begin{itemize}
  \item \Es{Krankentransportwagen} (\Es{KTW}), 
  \item \Es{Rettungstransportwagen}
(\Es{RTW}), 
\item \Es{Notarztwagen} (\Es{NAW}), 
\item \Es{Notarzteinsatzfahrzeug}
(\Es{NEWF}) und 
\item \Es{Notarzzthubschrauber} (\Es{NAH}) (in Deutschland:
Rettungstransporthubschrauber, RTH) 
\end{itemize}
eingebürgert.

\Layout{60}{Krankentransportwagen -- KTW}
{

Krankentransport.

Transport von erkrankten Patienten, die keine intensiven med. Massnahmen
benötigen. 

\blindtext

}{

}
{./icons/under-construction.png}
{}
{}

\Layout{60}{Rettungstransportwagen -- RTW}
{

Notfallrettung \& Überstellungen. 

Versorgung und Transport von erkrankenten Patienten und Notfallpatienten. 

Steigt ein Notarzt zu, sollte der RTW weitgehend gleichwertig zu einem NAW
sein. Zusammen mit dem NEF führt der RTW notarztpflichtige EInsätze durch. 

\blindtext
}{

}
{./images/IMG_1279.jpg}
{}
{}


\Layout{60}{Notarztwagen -- NAW}
{

Notfallrettung \& Überstellungen. 

Versorgung und Transport von Notfallpatienten. 

\blindtext

}{

}
{./icons/under-construction.png}
{}
{}


\Layout{60}{Notarzteinsatzfahrzeug -- NEF}
{
Notfallrettung \& Überstellungen. 

Versorgung von erkrankenten Notfallpatienten, \E{kein Transport}. Zubringer von
Notarzt. 

\blindtext
}{

}
{./images/nef_ma70_passat_nussdorferstrasse.jpg}
{}
{}


\Layout{60}{Notarzthubschrauber -- NAH}
{
Notfallrettung, Überstellungen, Bergungen. 

Versorgung und Transport von Notfallpatienten. 

\blindtext
}{

}
{./icons/under-construction.png}
{}
{}


\Layout{60}{Weitere Fahrzeugtypen}
{
PKW, Nk-KTW, B-KTW, SEW, MLS, ITW

\blindtext
}{

}
{./icons/under-construction.png}
{}
{}

% \begin{WidePar}

\Layout{57}{}
{}
{}
{\small
\begin{tabularx}{\linewidth}{XXXXXX}

&
\TabHeading{KTW}
&
\TabHeading{RTW}
&
\TabHeading{NAW}
&
\TabHeading{NEF}
&
\TabHeading{NAH}
\\\addlinespace\midrule\addlinespace
\TabSub{Name}
&
Krankentransportwagen
&
Rettungstransportwagen
&
Notarztwagen
&
Notarzteinsatzfahrzeug
&
Notarzthubschrauber
\\\addlinespace
\TabSub{Besatzung}
&
2 RS
&
1 NFS, 1-2 RS
&
1 NA, 1 NFS, 1-2 RS
\\\addlinespace
\TabSub{Funktion}
&	
Krankentransport.

Transport von erkrankten Patienten, die keine intensiven med. Massnahmen
benötigen. 
&
Notfallrettung \& Überstellungen. 

Versorgung und Transport von erkrankenten Patienten und Notfallpatienten. 
&
Notfallrettung \& Überstellungen. 

Versorgung und Transport von Notfallpatienten. 
&
Notfallrettung \& Überstellungen. 

Versorgung von erkrankenten Notfallpatienten, \E{kein Transport}. Zubringer von
Notarzt. 
&
Notfallrettung \& Überstellungen. 

Versorgung und Transport von Notfallpatienten. 
\\\addlinespace
\TabSub{Anmerkungen}

&

&
Steigt ein Notarzt zu, sollte der RTW weitgehend gleichwertig zu einem NAW sein. 
&

&

&
In Deutschland: "`RTH"'.
\\\addlinespace\bottomrule
\end{tabularx}
}
{Übersicht der wichtigsten Einstzmitteltypen in Österreich}
{}


% \end{WidePar}

\section{Regionales Rettungswesen}

\subsection{Rettungswesen in Wien}

\subsubsection{Organisationen}

\subsubsection{Organisation des Rettungsdienstes}

